\section{Introduction}
\label{sec:introduction}

Whole-body vibration (WBV) at \SIrange{4}{16}{\hertz} is a recognised
occupational hazard for vehicle operators.  ISO~2631-1~\cite{iso2631}
documents peak pelvic transmissibility at \SIrange{4}{8}{\hertz}, and
epidemiological studies report a consistent association between prolonged WBV
exposure and lower urinary tract symptoms---particularly increased urgency and
frequency---among truck drivers, bus operators, and forklift
operators~\cite{Seidel1986,griffin1990handbook}.  Despite this well-established
empirical record, no first-principles mechanical model of the urinary bladder
exists to explain \emph{why} whole-body vibration should preferentially excite
this organ, or to predict the frequency range at which the bladder wall is most
susceptible to vibratory loading.

The urinary bladder is, from a structural dynamics perspective, a remarkably
clean example of a fluid-filled viscoelastic shell.  The detrusor muscle forms
a roughly spherical wall of thickness \SIrange{1.5}{5}{\milli\metre}
(depending on fill state) enclosing urine---an incompressible fluid with
density and bulk modulus essentially identical to water.  The geometry, wall
stiffness, and internal pressure all vary systematically with fill volume,
making the bladder a natural candidate for parametric analysis across its
physiological operating range.

In a companion study~\cite{Mace2025browntone}, we developed an oblate
spheroidal shell model of the human abdominal cavity to investigate the
coupling between airborne infrasound and gastrointestinal resonance.  That
analysis revealed a coupling disparity of approximately $6.6 \times 10^4$
between mechanical and airborne excitation pathways---a result that explains
why whole-body vibration at \SIrange{4}{8}{\hertz} causes well-documented
gastrointestinal effects while airborne infrasound at comparable frequencies
does not.  The present work applies the same analytical framework to the
bladder, adapting the geometry, material properties, and boundary conditions to
this smaller, thinner-walled, and more compliant organ.

The mechanical properties of the bladder wall have been characterised by
several groups.  Ultrasound bladder vibrometry~\cite{Nenadic2013} measured
shear moduli of \SIrange{9.6}{48.7}{\kilo\pascal} across fill states from
\SI{187}{\milli\litre} to \SI{267}{\milli\litre}, demonstrating pronounced
strain-stiffening.  Barnes~\cite{Barnes2016} performed tensile testing of
excised bladder tissue, reporting Young's moduli from \SI{10}{\kilo\pascal} at
rest to several hundred kilopascals at high stretch ratios.  The wall exhibits
significant viscoelastic behaviour: stress relaxation, creep, and a loss
tangent estimated at 0.3--0.5 from dynamic
measurements~\cite{VanMastrigt1991}.  Wall thickness decreases from
approximately \SI{5}{\milli\metre} when empty to \SIrange{1.5}{3}{\milli\metre}
at capacity~\cite{Akkus2014}, consistent with a constant tissue volume
assumption.  Intravesical pressure follows the characteristic cystometric
curve: a compliant phase at low fill (\SIrange{5}{10}{\centi\metre}
water column) giving way to steep pressure rise above
\SI{300}{\milli\litre}~\cite{Griffiths1980}.

We apply these data to the fluid-filled shell framework.  The bladder is
modelled as a spherical shell (the near-spherical approximation is
appropriate here, unlike the distinctly oblate abdominal cavity) with
fill-dependent radius, wall thickness, elastic modulus, and internal pressure.
We compute the flexural mode frequencies across the physiological fill range
(\SIrange{50}{500}{\milli\litre}), identify the minimum-frequency fill state,
and compare airborne acoustic coupling with mechanical WBV coupling to the
bladder wall.

The remainder of this paper is organised as follows.
Section~\ref{sec:theory} presents the mathematical formulation, including the
fill-dependent parameter model and the modal analysis.
Section~\ref{sec:results} reports flexural mode frequencies across fill
volume, identifies the frequency minimum, and quantifies the coupling
disparity between mechanical and airborne excitation.
Section~\ref{sec:discussion} discusses clinical implications, limitations,
and comparison with the abdominal cavity model.
Section~\ref{sec:conclusion} draws conclusions.
